% =============================================================================
% CHAPTER 5: U-Net as Denoiser in Generative Models — EXPANDED
% =============================================================================
\section{U-Net as Denoiser in Generative Models}\index{U-Net}

\subsection{Original U-Net (Ronneberger et al., 2015)}\index{U-Net!Original}

\begin{definition}[U-Net]
A fully convolutional encoder-decoder with \textbf{symmetric skip connections} that concatenate encoder features with decoder features at corresponding resolution levels, forming a characteristic ``U'' shape.
\end{definition}

\textbf{Original purpose}: Biomedical image segmentation (cell tracking, neuron segmentation).\\
\textbf{Current role}: The \textbf{de facto backbone} for denoising in diffusion models (DDPM, Stable Diffusion, DALL-E 2, Imagen).

\subsection{U-Net Architecture — Layer by Layer}\index{U-Net!Architecture}

\textbf{Contracting Path (Encoder):}
\begin{center}\small
\begin{tabular}{clc}\toprule
\textbf{Level} & \textbf{Operations} & \textbf{Output Size} \\\midrule
Input & -- & $256 \times 256 \times 3$ \\
1 & $3{\times}3$ Conv $\times 2$, ReLU, MaxPool $2{\times}2$ & $128 \times 128 \times 64$ \\
2 & $3{\times}3$ Conv $\times 2$, ReLU, MaxPool $2{\times}2$ & $64 \times 64 \times 128$ \\
3 & $3{\times}3$ Conv $\times 2$, ReLU, MaxPool $2{\times}2$ & $32 \times 32 \times 256$ \\
4 & $3{\times}3$ Conv $\times 2$, ReLU, MaxPool $2{\times}2$ & $16 \times 16 \times 512$ \\
Bottleneck & $3{\times}3$ Conv $\times 2$, ReLU & $16 \times 16 \times 1024$ \\
\bottomrule
\end{tabular}
\end{center}

\textbf{Expanding Path (Decoder):}
\begin{center}\small
\begin{tabular}{clc}\toprule
\textbf{Level} & \textbf{Operations} & \textbf{Output Size} \\\midrule
4' & UpConv $2{\times}2$, \textbf{concat} Enc4, Conv$\times$2 & $32 \times 32 \times 512$ \\
3' & UpConv $2{\times}2$, \textbf{concat} Enc3, Conv$\times$2 & $64 \times 64 \times 256$ \\
2' & UpConv $2{\times}2$, \textbf{concat} Enc2, Conv$\times$2 & $128 \times 128 \times 128$ \\
1' & UpConv $2{\times}2$, \textbf{concat} Enc1, Conv$\times$2 & $256 \times 256 \times 64$ \\
Output & $1{\times}1$ Conv & $256 \times 256 \times C_{out}$ \\
\bottomrule
\end{tabular}
\end{center}

\subsubsection{Skip Connections — Why They're Critical}\index{U-Net!Skip Connections}
$$\mathbf{d}_\ell = [\text{UpConv}(\mathbf{d}_{\ell+1});\; \mathbf{e}_\ell] \quad \text{(channel concatenation)}$$

\begin{keybox}[Skip Connections — Complete Analysis]
\textbf{Without skip connections} (standard encoder-decoder):
\begin{itemize}
\item All information passes through the bottleneck ($16{\times}16$).
\item Fine spatial details (edges, textures) are lost in downsampling.
\item Decoder must \emph{hallucinate} fine details from compressed representation.
\item Result: blurry outputs, loss of spatial precision.
\end{itemize}

\textbf{With skip connections}:
\begin{itemize}
\item Encoder features at each level are directly available to decoder.
\item \textbf{Low-level features} (edges, gradients) from early encoder layers help decoder precisely localize structures.
\item \textbf{High-level semantics} from bottleneck guide \emph{what} to decode.
\item Combination: precise spatial reconstruction with semantic understanding.
\end{itemize}

\textbf{For denoising specifically}:
Noise exists across \textbf{all spatial frequencies}:
\begin{itemize}
\item \textbf{Low-frequency noise} (brightness variations): captured by deep features (bottleneck).
\item \textbf{High-frequency noise} (pixel-level Gaussian noise): captured by early encoder features.
\item Skip connections let the denoiser predict noise at ALL frequencies simultaneously.
\item Without skips: denoiser can remove low-freq noise but struggles with high-freq detail.
\end{itemize}

\textbf{Gradient flow benefit}: Direct paths from loss to early layers $\Rightarrow$ better training of shallow layers, especially important for deep U-Nets (50+ layers in some diffusion models).
\end{keybox}

\subsection{U-Net Modifications for Diffusion Models}\index{U-Net!Diffusion}

The standard U-Net is heavily modified for diffusion. Key additions:

\subsubsection{1. Time Conditioning}\index{U-Net!Time Conditioning}
The U-Net $\epsilon_\theta(\bx_t, t)$ must know the timestep $t$ to calibrate its denoising behavior. Different noise levels need different denoising strategies.

\textbf{Sinusoidal Time Embedding}\index{Sinusoidal Embedding} (same math as Transformer positional encoding):
\begin{equation}
\mathbf{t}_{emb}[2i] = \sin(t / 10000^{2i/d}), \quad \mathbf{t}_{emb}[2i+1] = \cos(t / 10000^{2i/d})
\end{equation}
Then project: $\mathbf{t}' = \text{SiLU}(\text{Linear}(\text{SiLU}(\text{Linear}(\mathbf{t}_{emb}))))$

\textbf{Injection via Adaptive Group Normalization (AdaGN)}\index{AdaGN}:
\begin{equation}
\boxed{\text{AdaGN}(\mathbf{h}, \mathbf{t}') = (1 + \gamma(\mathbf{t}')) \cdot \text{GroupNorm}(\mathbf{h}) + \beta(\mathbf{t}')}
\end{equation}
where $[\gamma, \beta] = \text{Linear}(\mathbf{t}')$ -- scale and bias from time embedding.

\begin{keybox}[Why AdaGN for Time Conditioning]
\begin{itemize}
\item \textbf{Why not just add $\mathbf{t}'$ to features?} Addition changes only the bias (shift). AdaGN provides both \textbf{scale} ($\gamma$) and \textbf{shift} ($\beta$), giving richer control.
\item \textbf{Why not concatenate $t$ as extra channel?} Time is a global property (same for all pixels). Concatenation treats it as spatial, which is wasteful.
\item \textbf{Why GroupNorm not BatchNorm?} GroupNorm works with small batch sizes (diffusion training often uses small effective batches). BN statistics are unstable with small batches.
\item \textbf{Intuition}: At high $t$ (lots of noise), $\gamma$ and $\beta$ should configure the network for ``coarse denoising.'' At low $t$ (little noise), configure for ``fine detail refinement.''
\end{itemize}
\end{keybox}

\subsubsection{2. Self-Attention Layers}\index{U-Net!Self-Attention}
Added at low-resolution feature maps (typically $16{\times}16$ and $8{\times}8$):
$$\text{SelfAttn}(\mathbf{Z}) = \softmax\!\left(\frac{(\mathbf{Z}\mathbf{W}_Q)(\mathbf{Z}\mathbf{W}_K)^\top}{\sqrt{d}}\right)(\mathbf{Z}\mathbf{W}_V) + \mathbf{Z}$$
where each spatial position is treated as a token.

\textbf{Why self-attention in denoiser?}
\begin{itemize}
\item Convolutions are local: $3{\times}3$ kernel sees 9 pixels. For generating symmetric faces, left-eye and right-eye need to coordinate across $\sim$100 pixels. Convolutions need $\sim$50 layers for this receptive field.
\item Self-attention: \textbf{global} connectivity in one layer. Left-eye features directly attend to right-eye features $\Rightarrow$ symmetry, coherence.
\item \textbf{Only at low-res}: At $16{\times}16$, sequence length = $256$ tokens. Attention cost: $256^2 \times d = $ manageable. At $256{\times}256$: $65536$ tokens $\Rightarrow$ $65536^2$ = 4 billion operations per layer — prohibitive!
\end{itemize}

\subsubsection{3. Cross-Attention for Text Conditioning}\index{U-Net!Cross-Attention}\index{Cross-Attention}
For text-to-image (Stable Diffusion), cross-attention layers inject text information:
\begin{equation}
\boxed{\text{CrossAttn}(\mathbf{Z}, \mathbf{c}) = \softmax\!\left(\frac{(\mathbf{Z}\mathbf{W}_Q)(\mathbf{c}\mathbf{W}_K)^\top}{\sqrt{d}}\right)(\mathbf{c}\mathbf{W}_V)}
\end{equation}
where $\mathbf{Z}$ = image features (\textbf{queries}), $\mathbf{c}$ = CLIP text embeddings (\textbf{keys and values}).

\begin{keybox}[Cross-Attention — How Text Controls Image]
\begin{itemize}
\item Each spatial position in image features ``queries'' the text tokens.
\item If spatial position $(x,y)$ corresponds to sky, it attends more to the word ``blue'' or ``cloudy.''
\item This creates \textbf{spatial-semantic binding}: specific image regions are guided by specific text words.
\item \textbf{Attention maps} reveal which text tokens influence which image regions. ``A cat sitting on a mat'': cat-region attends to ``cat'', mat-region attends to ``mat.'' 
\item This is why prompt structure matters: ``red car'' binds red to car; ``car red'' may not.
\end{itemize}
\end{keybox}

\subsubsection{4. ResNet Blocks in Diffusion U-Net}\index{U-Net!ResNet Blocks}
Each level uses \textbf{residual blocks} with time conditioning:
\begin{enumerate}
\item GroupNorm $\to$ SiLU $\to$ Conv3$\times$3
\item Add time embedding ($\mathbf{t}'$ projected to channel dimension)
\item GroupNorm $\to$ SiLU $\to$ Dropout $\to$ Conv3$\times$3
\item Residual connection ($1{\times}1$ conv if channels change)
\end{enumerate}

\textbf{Why residual blocks?}
\begin{itemize}
\item Deep U-Nets ($\sim$50 layers) need residual connections for gradient flow.
\item Each block learns a \textbf{refinement} (residual) rather than full transformation.
\item Better with time conditioning: the time embedding modulates the residual, adjusting the ``amount'' of denoising.
\end{itemize}

\subsection{Complete Diffusion U-Net Architecture}\index{U-Net!Full Architecture}

\textbf{Stable Diffusion v1.5 U-Net}: 860M parameters.

\begin{center}\small
\begin{tabular}{lcccc}\toprule
\textbf{Resolution} & \textbf{Channels} & \textbf{ResBlocks} & \textbf{Self-Attn} & \textbf{Cross-Attn} \\\midrule
$64 \times 64$ & 320 & 2 & No & Yes \\
$32 \times 32$ & 640 & 2 & Yes & Yes \\
$16 \times 16$ & 1280 & 2 & Yes & Yes \\
$8 \times 8$ (bottleneck) & 1280 & 1 & Yes & Yes \\
\bottomrule
\end{tabular}
\end{center}

\subsection{Latent Diffusion: U-Net in Compressed Space}\index{Latent Diffusion}\index{Stable Diffusion}

\textbf{Key insight}: Pixel-space diffusion on $512{\times}512$ is prohibitively expensive. Solution: denoise in \textbf{latent space}.

\begin{enumerate}
\item \textbf{VAE Encoder} $\mathcal{E}$: $512{\times}512{\times}3 \to 64{\times}64{\times}4$ (compression ratio: $48\times$).
\item \textbf{U-Net denoises latent}: $\epsilon_\theta(\bz_t, t, \mathbf{c})$ in $64{\times}64{\times}4$ space.
\item \textbf{VAE Decoder} $\mathcal{D}$: $64{\times}64{\times}4 \to 512{\times}512{\times}3$.
\end{enumerate}

\textbf{Compute savings}: Pixel-space attention at $64{\times}64$: $64^2 = 4096$ tokens. Pixel-space at $512{\times}512$: impossible (262144 tokens). Latent diffusion makes high-res generation tractable.

\subsection{ControlNet: Extending U-Net for Spatial Control}\index{ControlNet}
\textbf{Problem}: Text alone can't specify spatial layout (``cat on the LEFT'').

\textbf{ControlNet solution}: Clone the U-Net encoder. Feed conditioning signal (edge map, depth map, pose skeleton) through the clone. Add clone's outputs to the original U-Net's decoder via zero-initialized convolutions.

\textbf{Why ``zero convolution''?} Initially contributes nothing $\Rightarrow$ model starts from pre-trained weights unmodified. Gradually learns to incorporate spatial conditioning without destroying pre-trained quality.

\subsection{DiT: Replacing U-Net with Transformer}\index{DiT}
\textbf{Diffusion Transformer} (Peebles \& Xie, 2023): Replaces entire U-Net with a standard ViT.

\textbf{Architecture}: Patchify latent $\to$ add position embeddings $\to$ $N$ Transformer blocks with AdaLN (Adaptive LayerNorm for time/class conditioning) $\to$ unpatchify.

\textbf{Advantages over U-Net}:
\begin{itemize}
\item \textbf{Simpler}: No encoder-decoder, no skip connections, no upsampling.
\item \textbf{Scales better}: Clean power-law scaling with compute (Gflops $\propto$ quality).
\item \textbf{Leverages vision transformer advances}: ViT optimizations, FlashAttention, etc.
\end{itemize}
\textbf{Disadvantage}: At small scale, U-Net's inductive bias (multi-scale features) gives better results. DiT needs large compute to surpass U-Net.

\textbf{Trend}: SORA (video generation), Stable Diffusion 3, Flux all use Transformer-based architectures, replacing U-Net.

\subsection{When to Use U-Net / When NOT}\index{U-Net!When to Use}
\begin{center}\small
\begin{tabular}{p{0.45\linewidth}p{0.45\linewidth}}\toprule
\textbf{Use U-Net When} & \textbf{Avoid U-Net When} \\\midrule
Pixel/latent-level prediction needed & Simple classification (use encoder only) \\
Multi-scale features critical & Scalability is priority (use DiT) \\
Moderate resolution ($\leq 512^2$) & Very high resolution without cascading \\
Established diffusion pipeline & Cutting-edge scaling ($>$1B params, use Transformer) \\
Spatial correspondence in/out & Input/output have different structure \\
\bottomrule
\end{tabular}
\end{center}

\subsection{Exam Questions — U-Net (Expanded)}\index{U-Net!Exam Questions}

\begin{examq}[Q1: Architectural Analysis]
\textbf{Why are self-attention layers only added at low resolution ($\leq 32{\times}32$) in diffusion U-Nets?}

\textbf{Answer:} Self-attention has $O(N^2)$ complexity where $N$ = number of spatial positions. At $256{\times}256$: $N = 65536$, requiring $65536^2 \approx 4.3 \times 10^9$ operations per attention layer — computationally infeasible. At $16{\times}16$: $N = 256$, cost = $256^2 = 65536$ operations — manageable. Furthermore, at low resolution, each ``token'' has a large receptive field (each $16{\times}16$ position represents a $16{\times}16$ patch of the original), so attention at this level captures meaningful \textbf{global semantic relationships} (left eye $\leftrightarrow$ right eye, sky $\leftrightarrow$ ground). High-resolution convolutions handle local texture details efficiently.
\end{examq}

\begin{examq}[Q2: Design — Text Conditioning]
\textbf{Explain step-by-step how the prompt ``a red car on a green hill'' controls the U-Net denoising process in Stable Diffusion.}

\textbf{Answer:} (1) CLIP text encoder tokenizes and produces 77 embeddings of dim 768. Key tokens: ``red'', ``car'', ``green'', ``hill.'' (2) Noisy latent $\bz_t$ enters U-Net. At each level ($32^2$, $16^2$, $8^2$), image features are projected to queries $\mathbf{Q}$. (3) Text embeddings are projected to keys $\mathbf{K}$ and values $\mathbf{V}$. (4) Cross-attention: $\softmax(\mathbf{Q}\mathbf{K}^\top/\sqrt{d})\mathbf{V}$. Spatial positions representing the car region produce queries that have high dot product with ``red car'' keys $\Rightarrow$ attend to ``red car'' values. Hill region attends to ``green hill.'' (5) The attended text features are added to image features, biasing denoising toward producing red pixels in car-like regions and green pixels in hill-like regions. (6) This happens at \textbf{every denoising step} ($T$ times), iteratively sharpening the text-image alignment.
\end{examq}

\begin{examq}[Q3: Comparison]
\textbf{Compare U-Net vs DiT as diffusion backbones across 5 dimensions.}

\textbf{Answer:}
\begin{tabular}{lll}
\textbf{Dim} & \textbf{U-Net} & \textbf{DiT} \\
Architecture & Encoder-decoder + skips & Flat Transformer \\
Inductive bias & Multi-scale (conv) & Minimal \\
Scaling & Complex to scale & Clean power-law \\
Small compute & \textbf{Better} (bias helps) & Worse \\
Large compute & Good & \textbf{Better} (scaling law) \\
\end{tabular}
\end{examq}

\begin{examq}[Q4: Scenario]
\textbf{Your U-Net diffusion model generates images with good global structure but blurry textures. Diagnose and fix.}

\textbf{Answer:} \textbf{Diagnosis}: The model captures low-frequency content (global structure) but fails at high-frequency (texture). This suggests skip connections or high-resolution convolutions are weak. \textbf{Fixes}: (1) Add \textbf{more ResBlocks at high resolution} (Level 1, $128{\times}128$). Currently may have only 1; increase to 3. (2) Use \textbf{perceptual loss} (VGG feature matching) in addition to MSE noise prediction. (3) Apply \textbf{classifier-free guidance} with $w > 1$ at inference — amplifies conditional signal, sharpening details. (4) Check if \textbf{skip connections} are properly connected — missing skips would lose high-freq info. (5) Increase training resolution or use \textbf{progressive training} (train at $64^2$, then $256^2$).
\end{examq}
